% ============================================================================
%  TÀI LIỆU THAM KHẢO — xếp theo quy định: tiếng Việt trước (theo tên),
%  tiếng Anh sau (theo họ tác giả thứ nhất)
% ============================================================================
\renewcommand{\bibname}{TÀI LIỆU THAM KHẢO}
\begingroup
\small
\setstretch{1.2}
\begin{thebibliography}{99}
\addcontentsline{toc}{chapter}{Tài liệu tham khảo}
\setlength{\itemsep}{2pt}

\item[]\hspace*{-\labelwidth}\hspace*{-\labelsep}\textbf{Tiếng Việt}

\bibitem{dovanluu2000} Đỗ Văn Lưu, Phan Huy Khải (2000), \emph{Giải tích lồi}, Nhà xuất bản Khoa học và Kỹ thuật, Hà Nội.

\bibitem{vuhuutiep2018} Vũ Hữu Tiệp (2018), \emph{Machine Learning cơ bản}, Nhà xuất bản Khoa học và Kỹ thuật, Hà Nội.

\item[]\hspace*{-\labelwidth}\hspace*{-\labelsep}\textbf{Tiếng Anh}

\bibitem{arora2018} Arora R., Basu A., Mianjy P., Mukherjee A. (2018), “Understanding deep neural networks with rectified linear units”, in: \emph{International Conference on Learning Representations (ICLR 2018)}.

\bibitem{baesens2003} Baesens B., Van Gestel T., Viaene S., Stepanova M., Suykens J.A.K., Vanthienen J. (2003), “Benchmarking state-of-the-art classification algorithms for credit scoring”, \emph{Journal of the Operational Research Society}, 54(6), pp.~627--635.

\bibitem{bagirov2005} Bagirov A.M. (2005), “Max--min separability”, \emph{Optimization Methods and Software}, 20(2--3), pp.~277--296.

\bibitem{bennett2000} Bennett K.P., Bredensteiner E.J. (2000), “Duality and geometry in SVM classifiers”, in: \emph{Proceedings of the 17th International Conference on Machine Learning (ICML)}, pp.~57--64.

\bibitem{boser1992} Boser B.E., Guyon I.M., Vapnik V.N. (1992), “A training algorithm for optimal margin classifiers”, in: \emph{Proceedings of the 5th Annual Workshop on Computational Learning Theory (COLT)}, ACM, pp.~144--152.

\bibitem{boyd2004} Boyd S., Vandenberghe L. (2004), \emph{Convex Optimization}, Cambridge University Press, Cambridge.

\bibitem{breiman1993} Breiman L. (1993), “Hinging hyperplanes for regression, classification, and function approximation”, \emph{IEEE Transactions on Information Theory}, 39(3), pp.~999--1013.

\bibitem{canu2005} Canu S., Grandvalet Y., Guigue V., Rakotomamonjy A. (2005), \emph{SVM and Kernel Methods Matlab Toolbox}, Perception Systèmes et Information, INSA de Rouen, Rouen.

\bibitem{chang2011} Chang C.-C., Lin C.-J. (2011), “LIBSVM: a library for support vector machines”, \emph{ACM Transactions on Intelligent Systems and Technology}, 2(3), Article~27.

\bibitem{chua1977} Chua L.O., Kang S.M. (1977), “Section-wise piecewise-linear functions: canonical representation, properties, and applications”, \emph{Proceedings of the IEEE}, 65(6), pp.~915--929.

\bibitem{cortes1995} Cortes C., Vapnik V. (1995), “Support-vector networks”, \emph{Machine Learning}, 20(3), pp.~273--297.

\bibitem{cover1967} Cover T., Hart P. (1967), “Nearest neighbor pattern classification”, \emph{IEEE Transactions on Information Theory}, 13(1), pp.~21--27.

\bibitem{demsar2006} Demšar J. (2006), “Statistical comparisons of classifiers over multiple data sets”, \emph{Journal of Machine Learning Research}, 7, pp.~1--30.

\bibitem{do2017} Do T.-N., Poulet F. (2017), “Parallel learning of local SVM algorithms for classifying large datasets”, in: \emph{Transactions on Large-Scale Data- and Knowledge-Centered Systems XXXI}, Lecture Notes in Computer Science 10140, Springer, Berlin, Heidelberg.

\bibitem{fan2008} Fan R.-E., Chang K.-W., Hsieh C.-J., Wang X.-R., Lin C.-J. (2008), “LIBLINEAR: a library for large linear classification”, \emph{Journal of Machine Learning Research}, 9, pp.~1871--1874.

\bibitem{fawcett2006} Fawcett T. (2006), “An introduction to ROC analysis”, \emph{Pattern Recognition Letters}, 27(8), pp.~861--874.

\bibitem{freund1995} Freund Y., Schapire R.E. (1995), “A decision-theoretic generalization of on-line learning and an application to boosting”, in: \emph{Computational Learning Theory (EuroCOLT 1995)}, Lecture Notes in Computer Science 904, Springer, pp.~23--37.

\bibitem{friedman2001} Friedman J.H. (2001), “Greedy function approximation: a gradient boosting machine”, \emph{The Annals of Statistics}, 29(5), pp.~1189--1232.

\bibitem{goodfellow2013} Goodfellow I., Warde-Farley D., Mirza M., Courville A., Bengio Y. (2013), “Maxout networks”, in: \emph{Proceedings of the 30th International Conference on Machine Learning (ICML)}, pp.~1319--1327.

\bibitem{harris2020} Harris C.R. và cộng sự (2020), “Array programming with NumPy”, \emph{Nature}, 585(7825), pp.~357--362.

\bibitem{hastie1990} Hastie T., Tibshirani R. (1990), \emph{Generalized Additive Models}, Chapman \& Hall, London.

\bibitem{hsu2003} Hsu C.-W., Chang C.-C., Lin C.-J. (2003), \emph{A Practical Guide to Support Vector Classification}, Technical report, Department of Computer Science, National Taiwan University.

\bibitem{huanggb2006} Huang G.-B., Zhu Q.-Y., Siew C.-K. (2006), “Extreme learning machine: theory and applications”, \emph{Neurocomputing}, 70(1--3), pp.~489--501.

\bibitem{huang2013} Huang X., Mehrkanoon S., Suykens J.A.K. (2013), “Support vector machines with piecewise linear feature mapping”, \emph{Neurocomputing}, 117, pp.~118--127.

\bibitem{kelly_uci} Kelly M., Longjohn R., Nottingham K., \emph{The UCI Machine Learning Repository}, \texttt{https://archive.ics.uci.edu} (truy cập tháng 9/2026).

\bibitem{kostin2006} Kostin A. (2006), “A simple and fast multi-class piecewise linear pattern classifier”, \emph{Pattern Recognition}, 39(11), pp.~1949--1962.

\bibitem{lessmann2015} Lessmann S., Baesens B., Seow H.-V., Thomas L.C. (2015), “Benchmarking state-of-the-art classification algorithms for credit scoring: an update of research”, \emph{European Journal of Operational Research}, 247(1), pp.~124--136.

\bibitem{li2011} Li Y., Liu B., Yang X., Fu Y., Li H. (2011), “Multiconlitron: a general piecewise linear classifier”, \emph{IEEE Transactions on Neural Networks}, 22(2), pp.~276--289.

\bibitem{liu1989} Liu D.C., Nocedal J. (1989), “On the limited memory BFGS method for large scale optimization”, \emph{Mathematical Programming}, 45(1--3), pp.~503--528.

\bibitem{lundberg2017} Lundberg S.M., Lee S.-I. (2017), “A unified approach to interpreting model predictions”, in: \emph{Advances in Neural Information Processing Systems 30}, pp.~4765--4774.

\bibitem{maji2008} Maji S., Berg A.C., Malik J. (2008), “Classification using intersection kernel support vector machines is efficient”, in: \emph{IEEE Conference on Computer Vision and Pattern Recognition (CVPR)}, pp.~1--8.

\bibitem{mangasarian1968} Mangasarian O.L. (1968), “Multisurface method of pattern separation”, \emph{IEEE Transactions on Information Theory}, 14(6), pp.~801--807.

\bibitem{ovchinnikov2002} Ovchinnikov S. (2002), “Max-min representation of piecewise linear functions”, \emph{Beiträge zur Algebra und Geometrie}, 43(1), pp.~297--302.

\bibitem{pedregosa2011} Pedregosa F. và cộng sự (2011), “Scikit-learn: machine learning in Python”, \emph{Journal of Machine Learning Research}, 12, pp.~2825--2830.

\bibitem{platt1999} Platt J.C. (1999), “Fast training of support vector machines using sequential minimal optimization”, in: \emph{Advances in Kernel Methods: Support Vector Learning}, MIT Press, Cambridge, MA, pp.~185--208.

\bibitem{rahimi2007} Rahimi A., Recht B. (2007), “Random features for large-scale kernel machines”, in: \emph{Advances in Neural Information Processing Systems 20}, pp.~1177--1184.

\bibitem{scholkopf2002} Schölkopf B., Smola A.J. (2002), \emph{Learning with Kernels}, MIT Press, Cambridge, MA.

\bibitem{sklansky1980} Sklansky J., Michelotti L. (1980), “Locally trained piecewise linear classifiers”, \emph{IEEE Transactions on Pattern Analysis and Machine Intelligence}, 2(2), pp.~101--111.

\bibitem{stark1995} Stark P.B., Parker R.L. (1995), “Bounded-variable least-squares: an algorithm and applications”, \emph{Computational Statistics}, 10(2), pp.~129--141.

\bibitem{suykens1999ls} Suykens J.A.K., Vandewalle J. (1999), “Least squares support vector machine classifiers”, \emph{Neural Processing Letters}, 9(3), pp.~293--300.

\bibitem{suykens1999mlp} Suykens J.A.K., Vandewalle J. (1999), “Training multilayer perceptron classifiers based on a modified support vector method”, \emph{IEEE Transactions on Neural Networks}, 10(4), pp.~907--911.

\bibitem{suykens2002} Suykens J.A.K., Van Gestel T., De Brabanter J., De Moor B., Vandewalle J. (2002), \emph{Least Squares Support Vector Machines}, World Scientific, Singapore.

\bibitem{tao2022} Tao Q., Li L., Huang X., Xi X., Wang S., Suykens J.A.K. (2022), “Piecewise linear neural networks and deep learning”, \emph{Nature Reviews Methods Primers}, 2, Article~42.

\bibitem{tarela1999} Tarela J.M., Martínez M.V. (1999), “Region configurations for realizability of lattice piecewise-linear models”, \emph{Mathematical and Computer Modelling}, 30(11--12), pp.~17--27.

\bibitem{tenmoto1998} Tenmoto H., Kudo M., Shimbo M. (1998), “Piecewise linear classifiers with an appropriate number of hyperplanes”, \emph{Pattern Recognition}, 31(11), pp.~1627--1634.

\bibitem{thomas2002} Thomas L.C., Edelman D.B., Crook J.N. (2002), \emph{Credit Scoring and Its Applications}, SIAM, Philadelphia.

\bibitem{vanschoren2013} Vanschoren J., van Rijn J.N., Bischl B., Torgo L. (2013), “OpenML: networked science in machine learning”, \emph{SIGKDD Explorations}, 15(2), pp.~49--60.

\bibitem{vapnik1998} Vapnik V.N. (1998), \emph{Statistical Learning Theory}, Wiley, New York.

\bibitem{virtanen2020} Virtanen P. và cộng sự (2020), “SciPy 1.0: fundamental algorithms for scientific computing in Python”, \emph{Nature Methods}, 17(3), pp.~261--272.

\bibitem{wang2005} Wang S., Sun X. (2005), “Generalization of hinging hyperplanes”, \emph{IEEE Transactions on Information Theory}, 51(12), pp.~4425--4431.

\bibitem{webb2010} Webb D. (2010), \emph{Efficient Piecewise Linear Classifiers and Applications}, PhD thesis, University of Ballarat, Australia.

\bibitem{yeh2009} Yeh I.-C., Lien C.-H. (2009), “The comparisons of data mining techniques for the predictive accuracy of probability of default of credit card clients”, \emph{Expert Systems with Applications}, 36(2), pp.~2473--2480.

\end{thebibliography}
\endgroup
